\section{Introduction}

This report provides a diagnostic and rectification guide for the FESTO Processing Station. The aim of the lab work was to identify important failure situations, detect them using PLC logic, and provide clear lamp indications so that an operator or maintenance engineer can quickly find the possible cause.

The station performs a basic automated process: a part is detected in the socket, moved to the drilling position, drilled, checked by the testing cylinder, and then moved to the final position. The added diagnostic logic monitors this process using sensor feedback, timers, flags, and the three available indicator lamps.

The system uses the following diagnostic methods:

\begin{itemize}
    \item sensor inputs for material and position confirmation,
    \item timer-based detection for movement timeouts,
    \item logical checks for impossible sensor states,
    \item lamp indications for operator feedback,
    \item reset logic for fault recovery.
\end{itemize}

\section{Executive Summary}

Five fault categories were implemented for the FESTO Processing Station. They cover the main process, mechanical, pneumatic, electrical, and safety-related problems that can occur during operation.

\begin{longtable}{|p{0.8cm}|p{4.0cm}|p{4.0cm}|p{4.0cm}|}
\hline
No. & Fault Category & Root Cause Domain & Indication \\
\hline
1 & No part detected at start & Process / material & O1.0 blinking \\
\hline
2 & Test socket position fault & Mechanical / electrical / sensor & O1.0 + O1.1 blinking \\
\hline
3 & Drill cylinder fault & Pneumatic / mechanical / sensor & O1.0 + O1.2 blinking \\
\hline
4 & Testing cylinder fault & Pneumatic / mechanical / sensor & O1.1 + O1.2 blinking \\
\hline
5 & Stop or emergency stop condition & Operator / safety & O1.1 blinking \\
\hline
\end{longtable}

\section{Station Description}

The FESTO Processing Station consists of a part socket, rotating test socket/table, drilling mechanism, testing cylinder, and final part position. The automatic sequence starts when a part is placed in the socket and the Start button is pressed. The part is moved to the drilling position, drilled, checked by the testing cylinder, and transferred to the final position.

\begin{longtable}{|p{1.2cm}|p{4.2cm}|p{7.2cm}|}
\hline
Label & Station Part & Function \\
\hline
A & Part socket & Detects the inserted part \\
\hline
B & Drilling position & Holds the part for drilling \\
\hline
C & Drill & Performs the drilling operation \\
\hline
D & Testing cylinder & Checks the processed part \\
\hline
E & Final slot & Receives the processed part \\
\hline
\end{longtable}

\section{Allocation List}

\subsection{Inputs}

\begin{longtable}{|p{2.0cm}|p{10.8cm}|}
\hline
Address & Description \\
\hline
I0.0 & Part in socket \\
\hline
I0.1 & Test socket position sensor \\
\hline
I0.2 & Hold cylinder home position \\
\hline
I0.3 & Hold cylinder final position \\
\hline
I0.4 & Drill cylinder upper position \\
\hline
I0.5 & Drill cylinder lower position \\
\hline
I0.6 & Testing cylinder upper position \\
\hline
I0.7 & Testing cylinder lower position \\
\hline
I2.0 & Start button \\
\hline
I2.1 & Reset / clear diagnostics \\
\hline
I2.2 & Positioning button \\
\hline
I2.3 & Auto / manual control \\
\hline
I2.4 & Stop button \\
\hline
I2.7 & Emergency stop reset \\
\hline
\end{longtable}

\subsection{Outputs}

\begin{longtable}{|p{2.0cm}|p{10.8cm}|}
\hline
Address & Description \\
\hline
O0.1 & Drill engine \\
\hline
O0.2 & Test socket engine \\
\hline
O0.3 & Drill down command \\
\hline
O0.4 & Drill up command \\
\hline
O0.5 & Hold piston home command \\
\hline
O0.6 & Testing piston down command \\
\hline
O0.7 & Piston holder \\
\hline
O1.0 & Start lamp / diagnostic LED 1 \\
\hline
O1.1 & Reset lamp / diagnostic LED 2 \\
\hline
O1.2 & Positioning lamp / diagnostic LED 3 \\
\hline
\end{longtable}

\subsection{Added Diagnostic Flags and Timers}

\begin{longtable}{|p{2.4cm}|p{10.4cm}|}
\hline
Address & Description \\
\hline
F0.0 & Fault 1: no part detected \\
\hline
F0.1 & Fault 2: test socket position fault \\
\hline
F0.2 & Fault 3: drill cylinder fault \\
\hline
F0.3 & Fault 4: testing cylinder fault \\
\hline
F0.5 & Fault 5: stop or emergency stop condition \\
\hline
F0.6 & Diagnostic blink OFF phase \\
\hline
F0.7 & Diagnostic blink ON phase \\
\hline
T7 & Diagnostic LED ON timer \\
\hline
T8 & Diagnostic LED OFF timer \\
\hline
T9 & Movement fault timeout timer \\
\hline
\end{longtable}

\section{Reset Procedure}

The Reset button I2.1 clears active fault flags, blinking flags, diagnostic lamps, and diagnostic timers. After the fault is removed, the operator can reset the indication and return the station to normal operation.

\begin{verbatim}
IF I2.1 'Reset / clear diagnostics
THEN RESET F0.0
     RESET F0.1
     RESET F0.2
     RESET F0.3
     RESET F0.5
     RESET F0.6
     RESET F0.7
     RESET O1.0
     RESET O1.1
     RESET O1.2
     RESET T7
     RESET T8
     RESET T9
\end{verbatim}

\section{Detailed Fault Analysis}

\subsection{Fault 1: No Part Detected at Start}

\textbf{Root Cause Domain:} Process / material

\textbf{Severity:} High \quad \textbf{Detection Type:} Instantaneous

\subsubsection{Fault Condition}

This fault occurs when the operator presses Start I2.0, but the part sensor I0.0 does not detect a part in the socket. The cycle must not start because the station would attempt to move, drill, and test without material.

\subsubsection{Detection Mechanism}

\begin{longtable}{|p{3.6cm}|p{9.2cm}|}
\hline
Parameter & Details \\
\hline
Input checked & I0.0 part in socket \\
\hline
Trigger & I2.0 Start button pressed \\
\hline
Fault condition & I2.0 ON and I0.0 OFF \\
\hline
Diagnostic flag & F0.0 \\
\hline
Indicator & O1.0 blinking \\
\hline
\end{longtable}

\subsubsection{Possible Causes and Rectification}

\begin{itemize}
    \item Verify that a part is present in the socket.
    \item Check if the part is seated correctly.
    \item Inspect and clean sensor I0.0.
    \item Check the sensor LED, wiring, and connector.
    \item Replace or adjust the sensor if it does not switch.
    \item Press Reset I2.1 after correction and start again.
\end{itemize}

\subsubsection{Code Implementation}

\begin{verbatim}
IF I2.0
   AND I2.7
   AND N I0.0
THEN SET F0.0
     RESET F0.1
     RESET F0.2
     RESET F0.3
     RESET F0.5
\end{verbatim}

\begin{verbatim}
IF I2.0
   AND I2.7
   AND I0.0
THEN SET P2
\end{verbatim}

\subsection{Fault 2: Test Socket Position Fault}

\textbf{Root Cause Domain:} Mechanical / electrical / sensor

\textbf{Severity:} High \quad \textbf{Detection Type:} Timer-based

\subsubsection{Fault Condition}

This fault occurs when the test socket engine O0.2 is switched ON, but the position sensor I0.1 does not confirm the correct position within the timeout period. If this position is not confirmed, drilling or testing may be misaligned.

\subsubsection{Detection Mechanism}

\begin{longtable}{|p{3.6cm}|p{9.2cm}|}
\hline
Parameter & Details \\
\hline
Output command & O0.2 test socket engine \\
\hline
Sensor feedback & I0.1 test socket position sensor \\
\hline
Timer used & T9 movement fault timeout \\
\hline
Timeout value & 5 seconds \\
\hline
Fault condition & T9 finished and I0.1 still OFF \\
\hline
Diagnostic flag & F0.1 \\
\hline
Indicator & O1.0 + O1.1 blinking \\
\hline
\end{longtable}

\subsubsection{Possible Causes and Rectification}

\begin{itemize}
    \item Check if the table or test socket physically rotates.
    \item Inspect the table path for obstruction or mechanical debris.
    \item Check air supply or motor/actuator power depending on the mechanism.
    \item Verify that output O0.2 switches correctly.
    \item Inspect sensor I0.1 alignment and LED indication.
    \item Check PLC input wiring and connector for I0.1.
    \item Press Reset I2.1 after the fault is corrected.
\end{itemize}

\subsubsection{Code Implementation}

\begin{verbatim}
IF N T1
THEN SET O0.2
     SET T9 WITH 5s
\end{verbatim}

\begin{verbatim}
IF N T9
   AND N I0.1
THEN SET F0.1
     RESET O0.2
     JMP TO 200
\end{verbatim}

\subsection{Fault 3: Drill Cylinder Fault}

\textbf{Root Cause Domain:} Pneumatic / mechanical / sensor

\textbf{Severity:} Critical \quad \textbf{Detection Type:} Timer-based and logical conflict

\subsubsection{Fault Condition}

This fault occurs when the drill cylinder does not reach the commanded upper or lower position within the allowed time. It is also detected if both drill position sensors I0.4 and I0.5 are active at the same time, which is physically impossible.

\subsubsection{Detection Mechanism}

\begin{longtable}{|p{3.6cm}|p{9.2cm}|}
\hline
Parameter & Details \\
\hline
Down command & O0.3 drill down \\
\hline
Up command & O0.4 drill up \\
\hline
Lower sensor & I0.5 drill lower position \\
\hline
Upper sensor & I0.4 drill upper position \\
\hline
Timer used & T9 movement fault timeout \\
\hline
Timeout value & 3 seconds \\
\hline
Fault condition & Commanded position not reached, or I0.4 and I0.5 both ON \\
\hline
Diagnostic flag & F0.2 \\
\hline
Indicator & O1.0 + O1.2 blinking \\
\hline
\end{longtable}

\subsubsection{Possible Causes and Rectification}

\begin{itemize}
    \item Check pneumatic pressure and air supply to the drill cylinder.
    \item Inspect the drill cylinder for mechanical blockage or damage.
    \item Verify solenoid valve operation for O0.3 and O0.4.
    \item Check sensors I0.4 and I0.5 for alignment and LED state.
    \item Confirm that only one drill position sensor is active at a time.
    \item Inspect sensor wiring for loose wires or short circuits.
    \item Press Reset I2.1 after the fault is removed.
\end{itemize}

\subsubsection{Code Implementation}

\begin{verbatim}
IF O0.1
THEN SET O0.3
     SET T3 WITH 2s
     SET T4 WITH 1s
     SET T9 WITH 3s
\end{verbatim}

\begin{verbatim}
IF N T9
   AND N I0.5
THEN SET F0.2
     RESET O0.1
     RESET O0.3
     JMP TO 200
\end{verbatim}

\begin{verbatim}
IF N T9
   AND N I0.4
THEN SET F0.2
     RESET O0.1
     RESET O0.4
     JMP TO 200
\end{verbatim}

\begin{verbatim}
IF I0.4
   AND I0.5
THEN SET F0.2
\end{verbatim}

\subsection{Fault 4: Testing Cylinder Fault}

\textbf{Root Cause Domain:} Pneumatic / mechanical / sensor

\textbf{Severity:} High \quad \textbf{Detection Type:} Timer-based and logical conflict

\subsubsection{Fault Condition}

This fault occurs when the testing cylinder is commanded to move down using O0.6, but the lower position sensor I0.7 does not become active within the allowed time. A second fault condition is when upper sensor I0.6 and lower sensor I0.7 are active simultaneously.

\subsubsection{Detection Mechanism}

\begin{longtable}{|p{3.6cm}|p{9.2cm}|}
\hline
Parameter & Details \\
\hline
Output command & O0.6 testing piston down \\
\hline
Upper sensor & I0.6 testing cylinder upper position \\
\hline
Lower sensor & I0.7 testing cylinder lower position \\
\hline
Timer used & T9 movement fault timeout \\
\hline
Timeout value & 3 seconds \\
\hline
Fault condition & I0.7 not reached, or I0.6 and I0.7 both ON \\
\hline
Diagnostic flag & F0.3 \\
\hline
Indicator & O1.1 + O1.2 blinking \\
\hline
\end{longtable}

\subsubsection{Possible Causes and Rectification}

\begin{itemize}
    \item Check the testing cylinder for free movement.
    \item Inspect the testing probe for obstruction or incorrect part position.
    \item Check pneumatic pressure, tubing, and solenoid valve.
    \item Verify sensor I0.7 lower position detection.
    \item Verify sensor I0.6 upper position detection.
    \item Inspect wiring if both sensors are active together.
    \item Press Reset I2.1 after correction.
\end{itemize}

\subsubsection{Code Implementation}

\begin{verbatim}
IF N T2
THEN SET O0.6
     SET T4 WITH 1s
     SET T9 WITH 3s
\end{verbatim}

\begin{verbatim}
IF N T9
   AND N I0.7
THEN SET F0.3
     RESET O0.6
     JMP TO 200
\end{verbatim}

\begin{verbatim}
IF I0.6
   AND I0.7
THEN SET F0.3
\end{verbatim}

\subsection{Fault 5: Stop or Emergency Stop Condition}

\textbf{Root Cause Domain:} Operator / safety

\textbf{Severity:} Critical \quad \textbf{Detection Type:} Instantaneous

\subsubsection{Fault Condition}

This fault occurs when Stop I2.4 is pressed or when the emergency stop reset input I2.7 is not active. This condition requires the station to stop and remain stopped until the operator checks the reason.

\subsubsection{Detection Mechanism}

\begin{longtable}{|p{3.6cm}|p{9.2cm}|}
\hline
Parameter & Details \\
\hline
Stop input & I2.4 \\
\hline
Emergency reset input & I2.7 \\
\hline
Detection type & Instantaneous \\
\hline
Diagnostic flag & F0.5 \\
\hline
Indicator & O1.1 blinking \\
\hline
\end{longtable}

\subsubsection{Possible Causes and Rectification}

\begin{itemize}
    \item Identify why Stop or emergency stop was activated.
    \item Remove any obstruction or unsafe condition.
    \item Confirm that all motion has stopped.
    \item Reset the emergency stop circuit.
    \item Press Reset I2.1 to clear diagnostics.
    \item Restart only after the station is safe.
\end{itemize}

\subsubsection{Code Implementation}

\begin{verbatim}
IF I2.4
THEN SET F0.4
     SET F0.5
\end{verbatim}

\begin{verbatim}
IF N I2.7
THEN SET F0.5
     RESET O0.1
     RESET O0.2
     RESET O0.3
     RESET O0.4
     RESET O0.5
     RESET O0.6
     RESET O0.7
\end{verbatim}

\section{Lamp Indication Summary}

The three available lamps are reused as diagnostic indicators. The blinking pattern identifies the active fault and directs the operator to the first area to inspect.

\begin{longtable}{|p{0.8cm}|p{4.0cm}|p{3.4cm}|p{4.6cm}|}
\hline
No. & Fault Title & Lamp Pattern & First Inspection Area \\
\hline
1 & No part detected at start & O1.0 blinking & Part socket, I0.0 sensor \\
\hline
2 & Test socket position fault & O1.0 + O1.1 blinking & Table movement, I0.1 sensor \\
\hline
3 & Drill cylinder fault & O1.0 + O1.2 blinking & Drill cylinder, I0.4/I0.5 sensors \\
\hline
4 & Testing cylinder fault & O1.1 + O1.2 blinking & Testing cylinder, I0.6/I0.7 sensors \\
\hline
5 & Stop or emergency stop & O1.1 blinking & Stop button, emergency circuit \\
\hline
\end{longtable}

\section{Diagnostic Stop State}

When a diagnostic fault occurs during the automatic cycle, the program jumps to Step 200. This step removes the main motion commands before the operator resets or restarts the station.

\begin{verbatim}
STEP 200
IF NOP
THEN RESET O0.1
     RESET O0.2
     RESET O0.3
     RESET O0.4
     RESET O0.6
     RESET T9
     JMP TO 201
\end{verbatim}

\section{Testing and Validation}

The diagnostic functions were checked by creating each fault condition and observing the diagnostic flag, lamp indication, and program reaction. The expected result is that the correct lamp pattern appears and the unsafe output commands are removed.

\begin{longtable}{|p{1.0cm}|p{6.0cm}|p{5.8cm}|}
\hline
Test & Scenario & Expected Result \\
\hline
1 & Press Start with no part in socket & F0.0 active, O1.0 blinks, cycle does not start \\
\hline
2 & Block I0.1 during table movement & F0.1 active, O1.0 + O1.1 blink, movement stops \\
\hline
3 & Block I0.5 during drill down & F0.2 active, O1.0 + O1.2 blink, drill command stops \\
\hline
4 & Block I0.4 during drill return & F0.2 active, O1.0 + O1.2 blink, drill up command stops \\
\hline
5 & Block I0.7 during testing & F0.3 active, O1.1 + O1.2 blink, testing command stops \\
\hline
6 & Force I0.4 and I0.5 ON together & F0.2 active as drill sensor conflict \\
\hline
7 & Force I0.6 and I0.7 ON together & F0.3 active as testing sensor conflict \\
\hline
8 & Press Stop during automatic cycle & F0.5 active, O1.1 blinks, outputs stop \\
\hline
9 & Activate emergency stop condition & F0.5 active, O1.1 blinks, all main outputs reset \\
\hline
\end{longtable}

\section{Conclusion}

This report presents a structured diagnostic guide for the FESTO Processing Station. Five important faults were implemented: missing part, test socket position fault, drill cylinder fault, testing cylinder fault, and stop/emergency stop condition.

The diagnostic system uses PLC inputs, outputs, timers, flags, and lamp indications to identify faults and guide troubleshooting. Each fault includes the detection condition, possible causes, and rectification steps, so the report can be used as a practical manual for operators and maintenance engineers.

The final implementation improves fault detection, reduces troubleshooting time, and supports safer operation of the FESTO station.
